\chapter{Examen de Bloque B3}

\begin{exambrief}{Bloque B3: Trigonometria y Vectores}
Duracion sugerida: \(100\) minutos.
\begin{enumerate}
  \item Convierte \(240^\circ\) a radianes y determina exactamente las componentes del vector de
  longitud \(10\) con ese argumento.
  \item Resuelve en \([0,2\pi)\)
  \[
  2\sen^2 x-3\sen x+1=0.
  \]
  \item En un triangulo se conocen dos lados de \(\SI{11}{m}\) y \(\SI{14}{m}\), con angulo
  comprendido \(52^\circ\). Halla el tercer lado y el area.
  \item Dados \(\vec u=(3,2)\) y \(\vec v=(-1,4)\), calcula \(\vec u+\vec v\), el producto
  escalar y el angulo entre ambos vectores.
  \item Desde un punto de observacion se ve la cima de una torre con angulo \(30^\circ\). Tras
  avanzar \(\SI{20}{m}\) hacia la base, el angulo pasa a ser \(45^\circ\). Calcula la altura de
  la torre.
\end{enumerate}
\end{exambrief}

\begin{rubricbox}{Baremo orientativo B3}
\begin{itemize}
  \item Problema 1: \(2\) puntos.
  \item Problema 2: \(2\) puntos.
  \item Problema 3: \(2\) puntos.
  \item Problema 4: \(2\) puntos.
  \item Problema 5: \(2\) puntos.
\end{itemize}
\end{rubricbox}

\begin{shortanswer}{Bloque B3}
\[
240^\circ=\frac{4\pi}{3},
\qquad
(10\cos240^\circ,10\sen240^\circ)=(-5,-5\sqrt3).
\]
\[
x=\frac{\pi}{6},\ \frac{\pi}{2},\ \frac{5\pi}{6}.
\]
En el triangulo,
\[
c\approx \SI{11.29}{m},
\qquad
A\approx \SI{60.7}{m^2}.
\]
Para los vectores,
\[
\vec u+\vec v=(2,6),
\qquad
\vec u\cdot \vec v=5,
\qquad
\theta\approx 70.3^\circ.
\]
La altura de la torre es
\[
10(\sqrt3+1)\approx \SI{27.3}{m}.
\]
\end{shortanswer}

\begin{fullsolution}{Bloque B3}
1. Conversion y componentes:
\[
240^\circ=\frac{240\pi}{180}=\frac{4\pi}{3}.
\]
Como
\[
\cos240^\circ=-\frac12,
\qquad
\sen240^\circ=-\frac{\sqrt3}{2},
\]
las componentes del vector son
\[
(10\cos240^\circ,10\sen240^\circ)=(-5,-5\sqrt3).
\]

2. Ecuacion trigonometrica:
\[
2\sen^2x-3\sen x+1=0.
\]
Con \(u=\sen x\):
\[
2u^2-3u+1=(2u-1)(u-1)=0.
\]
Luego
\[
\sen x=\frac12 \quad \text{o}\quad \sen x=1.
\]
En \([0,2\pi)\),
\[
x=\frac{\pi}{6},\ \frac{5\pi}{6},\ \frac{\pi}{2}.
\]

3. Triangulo:
\[
c^2=11^2+14^2-2\cdot 11\cdot 14\cos52^\circ
=317-308\cos52^\circ.
\]
Como \(\cos52^\circ\approx 0.6157\),
\[
c^2\approx 317-189.64=127.36,
\qquad
c\approx \SI{11.29}{m}.
\]
El area es
\[
A=\frac12\cdot 11\cdot 14\sen52^\circ
=77\sen52^\circ\approx 77\cdot 0.7880\approx \SI{60.7}{m^2}.
\]

4. Vectores:
\[
\vec u+\vec v=(3,2)+(-1,4)=(2,6).
\]
El producto escalar es
\[
\vec u\cdot \vec v=3(-1)+2\cdot 4=5.
\]
Los modulos son
\[
\|\vec u\|=\sqrt{13},
\qquad
\|\vec v\|=\sqrt{17}.
\]
Entonces
\[
\cos\theta=\frac{5}{\sqrt{13}\sqrt{17}},
\qquad
\theta\approx 70.3^\circ.
\]

5. Torre:
Sea \(x\) la distancia desde el punto cercano a la base. Entonces el punto lejano esta a
\(x+20\). Por trigonometria,
\[
\tg45^\circ=\frac{h}{x}
\Longrightarrow h=x,
\]
\[
\tg30^\circ=\frac{h}{x+20}=\frac{x}{x+20}=\frac{1}{\sqrt3}.
\]
De aqui,
\[
\sqrt3 x=x+20
\Longrightarrow
x(\sqrt3-1)=20
\Longrightarrow
x=\frac{20}{\sqrt3-1}=10(\sqrt3+1).
\]
Como \(h=x\),
\[
h=10(\sqrt3+1)\approx \SI{27.3}{m}.
\]
\end{fullsolution}
